\section{Evaluation Pitfalls in Practice}
\label{sec:pitfalls}
% NEW (written for the TOMPECS version; not in the original paper)
Running the pipeline on production code exposed problems that benchmark runs had hidden, some of
them in the evaluation itself. Table~\ref{tab:pitfalls} collects them with the practice we adopted
in response; each is discussed where it arose, and the last two subsections of this section give
the tooling defects and the general contrast between benchmarks and applications.

% NEW: summary table. Every row restates a finding made elsewhere in the paper; no new measurement.
\begin{table}[t]
\centering
\footnotesize
\caption{Evaluation pitfalls observed in this study, where each is discussed, and the practice adopted.}
\label{tab:pitfalls}
\begin{tabular}{@{}p{0.27\textwidth}p{0.10\textwidth}p{0.52\textwidth}@{}}
\toprule
\textbf{Pitfall} & \textbf{Section} & \textbf{Practice adopted} \\
\midrule
Edits that delete the measured work while preserving output and the built-in check
  & \ref{sec:intent} & Semantic-preservation constraint; any speedup above $5\times$ is a correctness alarm and its diff is inspected \\
Built-in PASS/FAIL that gives no usable verdict on the unmodified program
  & \ref{sec:limits} & A candidate's verdict is informative only where it differs from the pristine one; explicit gates for applications \\
Drift between sessions on shared nodes (up to $49\%$ on identical code)
  & \ref{sec:timing} & Build once, interleave arms, wait for an idle GPU, log foreign GPU processes per repetition, require repetitions to agree \\
Application output with no built-in check, or not reproducible run to run
  & \ref{sec:realapps} & One gate per application, bit-exact or a statistical band; validate each gate by fault injection \\
Wall clock dominated by setup (AutoDock-GPU CUDA setup: $68$\,s)
  & \ref{sec:timing} & Use each application's own phase timer \\
Prompt contents differing from what was believed sent
  & \ref{sec:defects} & Defects fixed; earlier results reported as measured and not regenerated \\
Single samples and noisy pairs
  & \ref{sec:profabl}, \ref{sec:limits} & Score a pair a tie when its difference is inside its own spread; treat sub-$5\%$ single-cell effects as unresolved \\
Compilation context invisible to a single-file editor
  & \ref{sec:realapps} & Reported as the dominant real-code failure mode; carrying that context in the prompt is future work \\
\bottomrule
\end{tabular}
\end{table}

\subsection{Prompt-Fidelity Defects Exposed by Real Code}
\label{sec:defects}
Running production applications through the loop surfaced a class of problem that benchmark runs had
hidden for months: the prompt the model actually received was not the prompt we believed we were
sending. Table~\ref{tab:defects} lists the defects we found and fixed. They fall into two groups.
\emph{Summarization} defects corrupt the profiling evidence---a shared-memory footprint of
$7.43$\,Kbyte presented to the model as ``7 bytes'', SI \texttt{Kbyte} read as $1024$ bytes, and, most
seriously, a CSV reader that parsed every section of an Nsight Compute source page against the first
section's header and silently dropped rows of any other width. On one benchmark that left $36$ of $76$
stall rows in the prompt; re-running $56$ profiled inputs after the fix changed four of them, all by
\emph{gaining} rows. \emph{Loop} defects corrupt the repair channel: a build timeout reported to the
model as a fault in its own code, and an error filter that kept only lines matching
\texttt{error}$|$\texttt{fatal}$|$\texttt{undefined reference}, so that a linker's explanatory line
(\texttt{multiple definition of \ldots first defined here}) was discarded and the model saw only
\texttt{ld returned 1 exit status}.

One defect deserves separate mention because it is a design decision rather than an oversight. The
pipeline strips all comments from the source before prompting, to save tokens, and then replaces the
original file with the model's full-file answer---so every optimized file lost its license header and
all of its documentation. We initially, and wrongly, attributed the missing LGPL header of AutoDock-GPU to the
models; the prompts contained no occurrence of ``Copyright''. We now re-attach the original leading
comment block when it is a license and the candidate lacks it, but the general question stands: comment
stripping also discards institutional knowledge such as \texttt{// DO NOT UNDEFINE} directives that
exist precisely to stop the edit an optimizer is about to make.

\begin{table}[t]
\centering
\footnotesize
\setlength{\tabcolsep}{3pt}
\caption{Prompt-fidelity defects found while running production applications, and what each one
  corrupted. All were present during the benchmark study.}
\label{tab:defects}
\begin{tabular}{@{}p{0.30\columnwidth}p{0.63\columnwidth}@{}}
\toprule
\textbf{Defect} & \textbf{Effect} \\
\midrule
Shared memory read without unit & $7.43$\,Kbyte/block shown as ``7 bytes''; static and dynamic sizes
  summed across units \\
\texttt{Kbyte} treated as $1024$\,B & $2.4\%$ error in every stated size \\
Per-section CSV headers ignored & $25$ of $26$ sections dropped on one input ($36/76$ stall rows
  reached the model) \\
Stall reasons from a Python \texttt{set} & Prompt text not reproducible across processes \\
Kernel names matched verbatim & Namespaced and templated kernels fell back to the capped launch count \\
Comments stripped before prompting & Optimized files lost their license header and all documentation \\
Echoed line numbers not stripped & Uncompilable output when a model echoes the numbered source \\
Fixed $300$\,s build timeout & Real-application rebuilds reported to the model as its own fault \\
Linker cause lines filtered out & Repair prompt contained no actionable error \\
\bottomrule
\end{tabular}
\end{table}

\subsection{Benchmarks Versus Applications}
The benchmark results above understate how much prompt design matters, because single-kernel benchmarks
hide two of its effects. They cannot expose compilation-context failures, since each benchmark is a
self-contained translation unit; and they make profiling data look worthless
(Section~\ref{sec:profabl}), since the hot loop of a $200$-line benchmark is obvious from the source.
Both effects reappear on the production applications of Section~\ref{sec:realapps}, which is the
strongest argument we can offer against drawing conclusions about such systems from benchmark suites
alone.
